        \begin{keyinsight}[投资委员会结论]
        投资结论是谨慎参与、严格分层，而不是把钨、气体和氟化工当成同一批弹性股追高。当前行情真正交易的是半导体前道关键材料国产替代：WF6把钨粉、含氟反应、电子特气纯化和先进存储/AI算力扩产串成一条链。钨和制冷剂标的更像有当期业绩支撑的周期重估，电子特气更像订单和客户认证驱动的高弹性期权；下游晶圆厂和设备公司是需求锚，不应直接等同于上游利润兑现。
        \end{keyinsight}

        \section{先给组合结论：只做有证据的排序，不做主题平铺}
        这条主线最大的陷阱是把“产业重要性”误读成“所有相关股票都值得同样配置”。钨、制冷剂、电子特气都站在国产替代和资源安全的叙事里，但财报兑现路径完全不同：钨资源看价格和资源自给，制冷剂看配额和价差，WF6/NF3看高纯品级、晶圆厂认证、订单和毛利率。排序因此先按行动标签分层，再看最终目标价空间和证据质量；有估值空间但证据弱的标的不能上调为核心配置，有产业地位但价格已经显著透支的标的只能进入事件验证池。

        当前首选不是电子特气涨幅最大的股票，而是估值、业绩和证据更均衡的三美股份、翔鹭钨业、新宙邦、永和股份和章源钨业。三美股份的最终目标略高于现价，属于回撤配置；翔鹭钨业、新宙邦、永和股份和章源钨业的最终目标低于现价但差距仍可由下一季利润和价格信号验证，适合作为中性观察。巨化股份、多氟多和昊华科技开始进入高位风险区，原因不是产业逻辑消失，而是价格已经要求更强的二季度利润、订单或电子材料收入占比。

        \begin{exhibitbox}[核心排序与行动标签]
        \scriptsize
\resizebox{\exhibitwidth}{!}{
\begin{tabular}{cllllcll}
\toprule
\textbf{排序} & \textbf{代码} & \textbf{名称} & \textbf{现价} & \textbf{最终目标} & \textbf{空间} & \textbf{行动} & \textbf{失效条件} \\
\midrule
1 & 002842 & 翔鹭钨业 & 47.53 & 50.39 & 6.0\% & 回撤配置 & 钨价快速回落、矿端供给释放或下游硬质合金需求承压导致二季度利润环比下滑 \\
2 & 603379 & 三美股份 & 66.83 & 68.76 & 2.9\% & 回撤配置 & 制冷剂价格低于配额景气预期，或含氟材料/电子级产品认证低于预期 \\
3 & 300037 & 新宙邦 & 82.43 & 72.92 & -11.5\% & 中性观察 & 制冷剂价格低于配额景气预期，或含氟材料/电子级产品认证低于预期 \\
4 & 605020 & 永和股份 & 39.91 & 33.48 & -16.1\% & 中性观察 & 制冷剂价格低于配额景气预期，或含氟材料/电子级产品认证低于预期 \\
5 & 002378 & 章源钨业 & 33.78 & 27.26 & -19.3\% & 中性观察 & 钨价快速回落、矿端供给释放或下游硬质合金需求承压导致二季度利润环比下滑 \\
6 & 600160 & 巨化股份 & 50.72 & 40.38 & -20.4\% & 中性观察 & 制冷剂价格低于配额景气预期，或含氟材料/电子级产品认证低于预期 \\
7 & 600549 & 厦门钨业 & 81.03 & 57.47 & -29.1\% & 高位风险 & 钨价快速回落、矿端供给释放或下游硬质合金需求承压导致二季度利润环比下滑 \\
8 & 002407 & 多氟多 & 45.66 & 31.58 & -30.8\% & 高位风险 & 制冷剂价格低于配额景气预期，或含氟材料/电子级产品认证低于预期 \\
\bottomrule
\end{tabular}
}
\sourcenote{本报告模型，行情来自本地能力层实时抓取，财务为2026年一季报/2025年年报公开报表。}
\end{exhibitbox}

上一张表只放前八名，避免把投资委员会摘要写成全量数据库。但覆盖池共有18只，剩余标的同样需要给出行动标签和验证项；否则读者容易误以为没有进入前八名的电子特气和扩散标的没有结论。下面这张矩阵把其余10只放在同一口径下：价格和目标价用于判断当前赔率，行动标签用于决定研究行为，验证项用于决定后续能否从事件观察升级为基本面信用。

\begin{exhibitbox}[其余覆盖标的行动矩阵]
\scriptsize
\begin{tabularx}{\exhibitboxwidth}{L{0.55cm} L{0.85cm} L{1.1cm} L{1.15cm} L{1.15cm} L{0.8cm} L{1.2cm} L{0.9cm} X}
\toprule
\textbf{排序} & \textbf{代码} & \textbf{名称} & \textbf{主题} & \textbf{现价/目标} & \textbf{空间} & \textbf{行动} & \textbf{证据} & \textbf{下一步验证} \\
\midrule
9 & 600378 & 昊华科技 & 电子特气 & 60.13/36.96 & -38.5\% & 高位风险 & 高 & WF6收入占比、客户订单、价格和平台毛利 \\
10 & 300346 & 南大光电 & 电子特气 & 86.27/50.71 & -41.2\% & 高位风险 & 中 & 电子特气客户认证、订单、产品收入占比和毛利 \\
11 & 000657 & 中钨高新 & 钨 & 93.04/52.98 & -43.1\% & 高位风险 & 中 & 钨价、库存、二季度利润和现金流 \\
12 & 688146 & 中船特气 & 电子特气 & 339.00/174.66 & -48.5\% & 高位风险 & 高 & WF6/NF3订单、合同平均售价、收入确认和毛利 \\
13 & 688549 & 中巨芯-U & 电子特气 & 34.43/17.43 & -49.4\% & 高位风险 & 高 & WF6/NF3订单、合同平均售价、收入确认和毛利 \\
14 & 603505 & 金石资源 & 氟化工 & 18.85/8.38 & -55.6\% & 高位风险 & 中 & 制冷剂价差、电子材料客户和产品收入占比 \\
15 & 688268 & 华特气体 & 电子特气 & 240.90/102.57 & -57.4\% & 高位风险 & 高 & WF6收入占比、客户订单、价格和平台毛利 \\
16 & 002971 & 和远气体 & 电子特气 & 60.93/25.09 & -58.8\% & 高位风险 & 中 & 电子特气客户认证、订单、产品收入占比和毛利 \\
17 & 002549 & 凯美特气 & 电子特气 & 20.89/6.16 & -70.5\% & 高位风险 & 中 & 电子特气客户认证、订单、产品收入占比和毛利 \\
18 & 688106 & 金宏气体 & 电子特气 & 37.76/7.99 & -78.8\% & 高位风险 & 中 & 电子特气客户认证、订单、产品收入占比和毛利 \\
\bottomrule
\end{tabularx}
\sourcenote{与核心排序表同一估值模型、同一价格基准和同一证据门禁。未进入前八名不等于无研究价值；高位风险和事件验证标的需要订单、价格、客户或毛利证据后才可上修。}
\end{exhibitbox}

行动标签对应的是投资行为而不是情绪评价。优先跟踪表示当前价、最终目标和证据质量同时支持继续研究；回撤配置表示基本面逻辑成立，但需要更好价格或更明确财报确认；中性观察表示价格已经基本反映基本面预期，后续收益来自二/三季度超预期或事件验证；高位风险表示市场价格主要由远期情绪锚支撑，不能在缺少公告和利润验证时按低估处理。

\section{市场在为什么付钱：三类价值池不能混用倍数}
核心预期桥把股价背后的付费对象拆开。钨链的付费对象是战略资源、出口管制和钨价高位能否延续；制冷剂链的付费对象是配额约束、价差和现金流持续性；电子特气的付费对象是高纯WF6/NF3的国产认证、客户导入和潜在价格弹性。相同的“半导体材料”标签不能对应同一个市盈率，否则会把电子特气的远期期权错误套到钨加工或制冷剂现金流公司上。

\begin{exhibitbox}[核心标的市场预期桥]
\scriptsize
\resizebox{\exhibitwidth}{!}{
\begin{tabular}{llllllll}
\toprule
\textbf{代码} & \textbf{名称} & \textbf{现价} & \textbf{2026E收入} & \textbf{归母/每股收益} & \textbf{倍数} & \textbf{基本面锚} & \textbf{驱动} \\
\midrule
002842 & 翔鹭钨业 & 47.53 & 44.2亿元 & 9.6亿元/2.94 & 18.0x & 53.01 & 价格高位+资源稀缺+战略管制 \\
603379 & 三美股份 & 66.83 & 58.2亿元 & 21.1亿元/3.45 & 21.0x & 72.50 & 制冷剂价差+配额现金流 \\
300037 & 新宙邦 & 82.43 & 140.0亿元 & 20.0亿元/2.66 & 28.0x & 74.36 & 材料分部估值+电子级产品验证 \\
605020 & 永和股份 & 39.91 & 55.2亿元 & 7.5亿元/1.47 & 23.0x & 33.82 & 制冷剂价差+配额现金流 \\
002378 & 章源钨业 & 33.78 & 97.4亿元 & 14.1亿元/1.18 & 22.0x & 25.86 & 价格高位+资源稀缺+战略管制 \\
600160 & 巨化股份 & 50.72 & 250.7亿元 & 48.9亿元/1.81 & 22.0x & 39.83 & 制冷剂价差+配额现金流 \\
\bottomrule
\end{tabular}
}
\end{exhibitbox}

上表给出的基本面锚不是交易目标，而是财报和业务模型能够解释的价值中枢。三美股份、永和股份和巨化股份的基本面锚主要来自制冷剂价差现金流；翔鹭钨业和章源钨业的基本面锚来自钨价高位和资源稀缺；新宙邦更接近分部估值，需要电子氟化液和有机氟业务给出更清晰的利润贡献。只要下一季数据不能穿透到收入、毛利率和归母净利，主题热度就不能替代基本面锚。

高成长业绩精算给出的结论更严格：电子特气和电子材料第二曲线可以解释市场为什么愿意支付期权费，但还不能直接进入增长每股收益。只有订单量、平均售价、收入确认比例、分产品毛利和客户结构同步披露，产能或认证才可以从“事件验证”升级为“盈利信用”。

\begin{exhibitbox}[高成长业绩精算桥]
\scriptsize
\resizebox{\exhibitwidth}{!}{
\begin{tabular}{llllll}
\toprule
\textbf{代码} & \textbf{名称} & \textbf{增长段收入} & \textbf{销量/订单代理} & \textbf{增长每股收益} & \textbf{估值信用} \\
\midrule
688146 & 中船特气 & 收入/价格未披露 & WF6 2000吨/年；订单洽谈 & 无增长每股收益信用 & 订单+平均售价+毛利披露后再上修 \\
688549 & 中巨芯-U & 收入占比未披露 & WF6 600吨；无新增长单 & 无增长每股收益信用 & 订单+平均售价+毛利披露后再上修 \\
600378 & 昊华科技 & WF6约0.13\%收入占比 & 100吨/年项目口径 & 仅平台/分部估值观察 & 订单+平均售价+毛利披露后再上修 \\
688268 & 华特气体 & WF6占比较小 & 无合成产能；纯化供应 & 仅小权重观察 & 订单+平均售价+毛利披露后再上修 \\
\bottomrule
\end{tabular}
}
\end{exhibitbox}

\section{数据补齐到了哪一步：完整字段进估值，代理字段进监测，硬缺口阻断上修}
本报告已经把行情、官方市值、股本、财务质量、单季趋势、主营构成、产销量库存、客户认证、合同公告扫描、公开价格代理、券商样本、供应链关系、增长业绩和估值模型放进同一套数据室。真正仍缺的不是“有没有研究材料”，而是公司尚未公开披露的P0单品字段：合同平均售价、持续订单与收入确认比例、单品收入、客户份额和单品毛利。因此数据完整性审计的结论是有条件通过：完整字段可以进入估值，代理字段只能进入监测，硬缺口必须阻断增长每股收益上修。

\begin{exhibitbox}[数据完整性审计总览（核心字段）]
\scriptsize
\begin{tabularx}{\exhibitboxwidth}{L{2.1cm} L{1.8cm} X X}
\toprule
\textbf{数据域} & \textbf{覆盖/状态} & \textbf{模型纪律} & \textbf{仍缺/阻断} \\
\midrule
价格、股本、市值和流动性 & 18/18；完整 & 当前价和市值为估值主口径；流动性只约束拥挤度 & 交易所正式股本变动、自由流通股本和逐笔盘口仍需后续刷新 \\
基础财务、单季趋势和现金流质量 & 18/18；业绩披露日历18/18；业绩预告/快报命中16/18；完整 & 可用于基准业务估值和披露窗口管理；预约日期、业绩预告或快报均不自动证明单品增长 & 产品级季度收入、项目级回款、实际半年报结果、券商一致预期和单品级业绩拆分未接入 \\
主营构成、官方分部和产品毛利 & 18/18；核心电子特气官方分部4家公司/9个分部；价格与毛利入模门禁18/18；合同平均售价入模0/18；单品毛利入模0/18；有条件完整 & 大类分部可作上限；合同平均售价和单品毛利缺失前不能给增长每股收益 & WF6/NF3/含氟电子材料/高纯钨材单品收入、单品毛利和客户份额 \\
产销量、库存、产能利用和项目建设 & 18/18；有条件完整 & 产能和销量只能进监测与直接性；不能直接折成增长收入 & 单品订单、合同平均售价、收入确认比例和单品毛利 \\
客户链、认证、前五客户和公司沟通 & 18/18；高成长客户认证8家公司；客户侧验证锚点8个；下游客户公开年报8/8份；有条件完整 & 客户认证、供应商侧客户披露和下游客户年报产品词命中不等于客户侧采购确认；匿名客户不做晶圆厂推断 & 单品客户份额、客户侧供应商名录、客户侧采购订单、订单金额和合同价格 \\
\bottomrule
\end{tabularx}
\sourcenote{完整审计见 data/data\_completeness\_dashboard.json/md 与 analysis/data\_completeness\_dashboard.md。}
\end{exhibitbox}

\begin{exhibitbox}[数据完整性审计总览（代理与门禁）]
\scriptsize
\begin{tabularx}{\exhibitboxwidth}{L{2.1cm} L{1.8cm} X X}
\toprule
\textbf{数据域} & \textbf{覆盖/状态} & \textbf{模型纪律} & \textbf{仍缺/阻断} \\
\midrule
合同、订单、收入确认和履约代理 & 公告扫描18/18；正文解析25条；主题订单金额样本2条；合同金额经济性约束3条；订单耐久性门禁18/18，持续订单增长每股收益0/18；单品边界5只/9行；问询/审核回复命中3/18；部分补证 & 中船历史订单只提高事件验证；订单耐久性门禁显示0/18可进入持续订单增长每股收益；问询回复只有披露可量化单品字段才改变增长每股收益信用 & 持续订单、合同平均售价、交付节奏、客户份额和单品毛利 \\
公开价格、下游需求、AI平台和政策锚 & 价格代理9项；公司级HFC配额和全国份额4只；下游需求5个；AI平台7条；政策/制程/产品能力均已归档；代理完整 & HFC配额和份额可支持制冷剂基准现金流及公司相对地位；其他价格/需求锚只作代理和监测，不替代公司收入或利润 & 公司合同价格、客户采购量、电子材料订单/价差、收入确认和单品毛利 \\
高成长业绩、单品转化和隐含增长 & 高成长标的14/14；增长每股收益允许入模0/14；证据阶段：事件验证池3只；弱观察池10只；强观察池5只；阻断入模 & P0未关闭前，高成长段不能进入基本面目标价上修 & 4类P0硬缺口仍未关闭：合同价格、订单确认、单品收入和单品毛利 \\
外部券商、来源治理和字段矩阵 & 公开研报样本18/18；公开PDF全文18篇；字段矩阵65组；来源抓取失败0条；有条件完整 & 公开摘要和公开PDF全文只能作语境与复核锚；不替代本报告估值门禁 & 付费一致预期、完整券商模型表、客户侧文件和更多原始调研纪要 \\
\bottomrule
\end{tabularx}
\sourcenote{有条件完整不等于P0关闭；合同平均售价、单品收入、客户份额和单品毛利未披露前，高成长段不得进入基本面目标价上修。}
\end{exhibitbox}

\section{情绪锚只解释价格，不替代基本面}
主题行情里，市场经常先给稀缺资产一个情绪锚，再等公告和财报补证据。本模型保留市场锚，是为了避免机械地把强势股全部判成卖出；同时又把基本面锚、市场锚和券商锚拆开，防止用成交热度替代估值结论。权重中的基本面/市场/外部分别代表基本面、市场情绪和外部券商或公告锚，电子特气的市场权重更高，但证伪也更快。

\begin{exhibitbox}[市场隐含情绪锚和最终权重]
\scriptsize
\resizebox{\exhibitwidth}{!}{
\begin{tabular}{llllllll}
\toprule
\textbf{代码} & \textbf{名称} & \textbf{内在锚} & \textbf{市场锚} & \textbf{券商锚} & \textbf{权重（基本面/市场/外部）} & \textbf{最终目标} & \textbf{行动逻辑} \\
\midrule
002842 & 翔鹭钨业 & 53.01 & 42.54 & 未披露 & 75\%/25\%/0\% & 50.39 & 基本面和估值仍可跟踪，等待价格或财报确认 \\
603379 & 三美股份 & 72.50 & 53.80 & 未披露 & 80\%/20\%/0\% & 68.76 & 基本面和估值仍可跟踪，等待价格或财报确认 \\
300037 & 新宙邦 & 74.36 & 66.36 & 未披露 & 82\%/18\%/0\% & 72.92 & 中性观察，等待下一季业绩和价格信号 \\
605020 & 永和股份 & 33.82 & 32.13 & 未披露 & 80\%/20\%/0\% & 33.48 & 中性观察，等待下一季业绩和价格信号 \\
002378 & 章源钨业 & 25.86 & 30.23 & 未披露 & 68\%/32\%/0\% & 27.26 & 中性观察，等待下一季业绩和价格信号 \\
600160 & 巨化股份 & 39.83 & 42.35 & 未披露 & 78\%/22\%/0\% & 40.38 & 中性观察，等待下一季业绩和价格信号 \\
\bottomrule
\end{tabular}
}
\end{exhibitbox}

对投资委员会来说，情绪锚的用法很直接：如果市场锚高、但基本面锚低，股票只能作为事件验证；如果基本面锚与市场锚接近，回撤后才有配置意义。三美股份和翔鹭钨业的基本面锚与当前价相对接近，安全边际问题可以通过价格和二季度利润解决；昊华科技、电子特气扩散股以及高拥挤标的则需要订单、价格或收入占比证明，否则应下调市场权重。

\section{下一季是验证点：没有二季度桥，主题就无法升级为业绩}
这条链后续不缺故事，缺的是能把故事转成财报的指标。二季度最低门槛给的是研究上的保底验证线：收入、归母净利、价差、钨价或客户认证如果达不到，最终目标要向基本面锚回归；如果达到并伴随公告级订单或价格证据，才可以把市场锚保留到下一轮模型更新。

\begin{exhibitbox}[下一季度最低验证门槛]
\scriptsize
\resizebox{\exhibitwidth}{!}{
\begin{tabular}{llll}
\toprule
\textbf{代码} & \textbf{名称} & \textbf{二季度门槛} & \textbf{催化} \\
\midrule
002842 & 翔鹭钨业 & 二季度收入约11.0亿元、归母约2.4亿元且钨价不回落 & 钨矿/钨粉价格维持高位、出口许可偏紧、光伏钨丝与军工高温材料需求兑现 \\
603379 & 三美股份 & 二季度收入约14.6亿元、归母约5.3亿元且制冷剂/电子材料价差维持 & R32/R125/R134a配额约束和价差修复、电子级HF/氟化液客户验证推进 \\
300037 & 新宙邦 & 二季度收入约35.0亿元、归母约5.0亿元且制冷剂/电子材料价差维持 & R32/R125/R134a配额约束和价差修复、电子级HF/氟化液客户验证推进 \\
605020 & 永和股份 & 二季度收入约13.8亿元、归母约1.9亿元且制冷剂/电子材料价差维持 & R32/R125/R134a配额约束和价差修复、电子级HF/氟化液客户验证推进 \\
002378 & 章源钨业 & 二季度收入约24.4亿元、归母约3.5亿元且钨价不回落 & 钨矿/钨粉价格维持高位、出口许可偏紧、光伏钨丝与军工高温材料需求兑现 \\
600160 & 巨化股份 & 二季度收入约62.7亿元、归母约12.2亿元且制冷剂/电子材料价差维持 & R32/R125/R134a配额约束和价差修复、电子级HF/氟化液客户验证推进 \\
\bottomrule
\end{tabular}
}
\end{exhibitbox}

验证门槛必须绑定披露窗口，否则“下一季验证”会变成口号。东方财富业绩预约披露日历显示，中船特气2026年半年报预约披露日为2026-07-18；预约披露日只触发模型刷新，不替代实际财报结果、订单、合同平均售价、收入确认比例、客户份额、单品收入或单品毛利。所有标的在预约披露日后都要重新核验收入、利润、毛利、现金流和公告级订单，未披露硬字段时仍按期权或观察信用处理。

\begin{exhibitbox}[业绩披露日历与验证窗口]
\scriptsize
\begin{tabularx}{\exhibitboxwidth}{L{0.85cm} L{1.15cm} L{1.25cm} L{1.35cm} L{0.85cm} X}
\toprule
\textbf{代码} & \textbf{名称} & \textbf{下一报告} & \textbf{预约披露} & \textbf{状态} & \textbf{验证重点} \\
\midrule
002842 & 翔鹭钨业 & 2026年 一季报 & 2026-04-29 & 已披露 & 核验钨价、库存、二季度利润、现金流和高纯钨材客户/订单。 \\
603379 & 三美股份 & 2026年 半年报 & 2026-08-25 & 待披露 & 核验制冷剂价差、HFC配额兑现、电子材料客户/订单/价差。 \\
300037 & 新宙邦 & 2026年 一季报 & 2026-04-28 & 已披露 & 核验制冷剂价差、HFC配额兑现、电子材料客户/订单/价差。 \\
605020 & 永和股份 & 2026年 半年报 & 2026-08-17 & 待披露 & 核验制冷剂价差、HFC配额兑现、电子材料客户/订单/价差。 \\
002378 & 章源钨业 & 2026年 一季报 & 2026-04-28 & 已披露 & 核验钨价、库存、二季度利润、现金流和高纯钨材客户/订单。 \\
600160 & 巨化股份 & 2026年 半年报 & 2026-08-26 & 待披露 & 核验制冷剂价差、HFC配额兑现、电子材料客户/订单/价差。 \\
688146 & 中船特气 & 2026年 半年报 & 2026-07-18 & 待披露 & 核验WF6/NF3订单金额、合同均价、收入确认、客户份额和单品毛利。 \\
\bottomrule
\end{tabularx}
\sourcenote{东方财富业绩预约披露日历接口；完整18只标的见 data/earnings\_disclosure\_calendar.json/md。预约披露日仅定义复核窗口，不替代实际财报结果、订单、合同平均售价、收入确认比例、客户份额、单品收入或单品毛利。}
\end{exhibitbox}

业绩预告和快报证据进一步回答：在半年报正式披露前，公司是否已经用公告形式给出利润预警、快报或同比变化线索。本报告只把它作为公司层业绩验证入口；即使公告显示利润改善，也必须继续拆到产品、订单、合同平均售价、收入确认比例、客户份额和单品毛利，才可以改变高成长段的增长每股收益信用。

\begin{exhibitbox}[业绩预告/快报公告证据]
\scriptsize
\begin{tabularx}{\exhibitboxwidth}{L{0.85cm} L{1.1cm} L{1.25cm} L{1.05cm} X X L{1.45cm}}
\toprule
\textbf{代码} & \textbf{名称} & \textbf{日期} & \textbf{类型} & \textbf{标题} & \textbf{业绩字段状态} & \textbf{估值处理} \\
\midrule
688106 & 金宏气体 & 2026-04-29 & 业绩预告 & 金宏气体:2026年第一季度业绩预告的自愿性披露公告 & 正文含利润/同比变动和原因解释，可用于业绩验证但不关闭单品P0字段 & 业绩验证；不替代单品硬字段 \\
600378 & 昊华科技 & 2026-04-24 & 业绩预增 & 昊华科技:昊华科技2026年第一季度业绩预增公告 & 正文含利润/同比变动和原因解释，可用于业绩验证但不关闭单品P0字段 & 业绩验证；不替代单品硬字段 \\
300037 & 新宙邦 & 2026-04-10 & 业绩预告 & 新宙邦:(2026-033)2026年第一季度业绩预告 & 正文含利润或同比变动区间，可用于业绩验证但需核对口径 & 业绩验证；不替代单品硬字段 \\
000657 & 中钨高新 & 2026-04-10 & 业绩快报 & 中钨高新:中钨高新材料股份有限公司2025年度业绩快报 & 正文含利润/同比变动和原因解释，可用于业绩验证但不关闭单品P0字段 & 业绩验证；不替代单品硬字段 \\
000657 & 中钨高新 & 2026-04-10 & 业绩预增 & 中钨高新:中钨高新材料股份有限公司2026年一季度业绩预增的公告 & 正文含利润/同比变动和原因解释，可用于业绩验证但不关闭单品P0字段 & 业绩验证；不替代单品硬字段 \\
002971 & 和远气体 & 2026-04-09 & 业绩快报 & 和远气体:2025年度业绩快报 & 正文含利润/同比变动和原因解释，可用于业绩验证但不关闭单品P0字段 & 业绩验证；不替代单品硬字段 \\
002378 & 章源钨业 & 2026-03-14 & 业绩快报 & 章源钨业:2025年度业绩快报 & 正文含利润/同比变动和原因解释，可用于业绩验证但不关闭单品P0字段 & 业绩验证；不替代单品硬字段 \\
688268 & 华特气体 & 2026-02-28 & 业绩快报 & 华特气体:广东华特气体股份有限公司2025年度业绩快报公告 & 正文含利润/同比变动和原因解释，可用于业绩验证但不关闭单品P0字段 & 业绩验证；不替代单品硬字段 \\
\bottomrule
\end{tabularx}
\sourcenote{东方财富公告索引与正文接口；扫描业绩预告、业绩快报、业绩预增/预减/预亏/预盈公告，排除业绩说明会、业绩承诺和专项审核类公告。完整18只标的见 data/performance\_forecast\_evidence.json/md。}
\end{exhibitbox}

行动上分三条路径。第一，回撤配置路径关注三美股份、翔鹭钨业等基本面锚相对扎实的标的，等待价格给出安全边际。第二，中性观察路径关注新宙邦、永和股份、章源钨业，重点看二季度利润和产品结构能否抬高基本面锚。第三，事件验证路径关注WF6/NF3电子特气和扩散气体标的，只在公告、客户认证、价格或订单出现后提高权重，不能把题材热度本身当成业绩兑现。
